\section{CEO 视角：迷茫期、心力与 All In}

本期另一个重要层面，是创业者如何管理探索期。沐言智语融资顺利，但张月光反复说，钱不是最重要的，团队心力和自己的心力更重要。没有结果的探索会消耗团队，如果创始人不断换方向，又不管理预期，团队会被搞崩。

\subsection{迷茫期如何沟通}

他描述自己在迷茫期的原则：实事求是，不把没有信心的想法包装成确定方向；告诉团队现在还没有 all in；允许别人提出更好的想法；用小成本试错控制团队心力。这个做法和产品探索方法一致：不确定时控制变量，确定后才集中资源。

\begin{knowledgebox}{探索期的组织管理}
AI 应用创业很容易陷入“每个方向都像机会”的状态。探索期要保护现金，更要保护团队心力；all in 之前，要明确告诉团队这只是试错，而不是最终战役。
\end{knowledgebox}

\subsection{为什么 Dokie 可以 All In}

张月光说，当他想到 Dokie 背后那件事——做一个突破自己能力边界的 AI 朋友——即使把公司的钱输在这件事上，他也可以接受。这个标准很硬：不是看单个 PPT 产品能不能爆，而是看背后的大命题是否值得公司押注。

\begin{importantbox}{All In 的判断}
All in 不应来自焦虑、融资压力或追热点，而应来自创始人对大命题的确信：即使失败，也知道为什么值得输在这里。
\end{importantbox}

\subsection{本章小结}

AI 应用公司不仅要找产品，也要管理探索期。张月光对 Dokie 的 all in，真正押注的是“突破能力边界的 AI 朋友”，而不是 PPT 这个单点工具。

